% Generated draft: requires mathematical and source review.
\documentclass[11pt]{article}
\usepackage[margin=25mm]{geometry}
\usepackage{fontspec}
\setmainfont{Latin Modern Roman}
\usepackage{amsmath,amssymb,mathtools,unicode-math}
\setmathfont{Latin Modern Math}
\usepackage{enumitem,needspace,xurl,hyperref}
\setlength{\parindent}{0pt}
\setlength{\parskip}{5pt}
\setlength{\emergencystretch}{3em}
\newcommand{\N}{\mathbb{N}}
\newcommand{\Z}{\mathbb{Z}}
\newcommand{\Q}{\mathbb{Q}}
\newcommand{\R}{\mathbb{R}}
\newcommand{\C}{\mathbb{C}}
\newcommand{\F}{\mathbb{F}}
\newcommand{\A}{\mathbb{A}}
\newcommand{\PP}{\mathbb{P}}
\newcommand{\E}{\mathbb{E}}
\newcommand{\eps}{\varepsilon}
\newcommand{\dd}{\,\mathrm{d}}
\begin{document}
\section*{UnsolvedMath ID: 800622}
\textbf{Source problem code:} catalogue entry
\subsection{Definitions and mathematical statement}
For a topological space \(T\), let \(\mathcal O(T)\) be the family of open covers of \(T\). An open \(\omega\)-cover of \(T\) is a family \(\mathcal U\) of open subsets of \(T\) such that \(T\notin\mathcal U\) and every finite subset of \(T\) is contained in some member of \(\mathcal U\). Write \(\Omega(T)\) for the family of these covers. Under the standard selection-principle convention, \(T\) satisfies \(\mathsf U_{\mathrm{fin}}(\mathcal O,\Omega)\) if, for every sequence \((\mathcal U_n)_{n\in\mathbb N}\) of open covers of \(T\), each having no finite subcover, there are finite subfamilies \(\mathcal F_n\subseteq\mathcal U_n\) such that
\[
\left\{\bigcup\mathcal F_n:n\in\mathbb N\right\}\in\Omega(T).
\]
For functions \(f,g\in\mathbb N^{\mathbb N}\), define eventual domination by
\[
g\leq^* f\quad\Longleftrightarrow\quad
\{n\in\mathbb N:g(n)>f(n)\}\text{ is finite}.
\]
For \(k\geq1\), the pointwise maximum of \(f_1,\ldots,f_k\in\mathbb N^{\mathbb N}\) is the function \(n\mapsto\max_{1\leq i\leq k}f_i(n)\).

The supplied Problem 51 assumes that \(X\) satisfies \(\mathsf U_{\mathrm{fin}}(\mathcal O,\Omega)\) and asks whether the set displayed in the source as
\[
G=\{g\in\mathbb N:(\exists k)(\exists f_1,\ldots,f_k\in Y)\ g\leq^*\max\{f_1,\ldots,f_k\}\}
\]
also satisfies \(\mathsf U_{\mathrm{fin}}(\mathcal O,\Omega)\).

This formulation is incomplete: \(Y\) is undefined, its relationship to \(X\) is unspecified, and no topology on either space is supplied. Moreover, the membership condition \(g\in\mathbb N\) does not provide a function to which the usual eventual-domination relation applies. If the intended membership condition is \(g\in\mathbb N^{\mathbb N}\), if \(Y\subseteq\mathbb N^{\mathbb N}\), and if only nonempty finite maxima are intended, then the corresponding well-defined set is
\[
G_Y=\left\{g\in\mathbb N^{\mathbb N}:\exists k\in\mathbb N,\ k\geq1,\ \exists f_1,\ldots,f_k\in Y,\ g\leq^*\max_{1\leq i\leq k}f_i\right\}.
\]
These are conditional interpretations, not supplied hypotheses. The record does not determine which relationship between \(X\) and \(Y\), or which topology on \(G_Y\), is intended, so it does not yet specify a precise implication to be assessed.
\subsection{Short English statement}
The source asks whether a set obtained by closing \(Y\) downward under eventual domination by finite pointwise maxima inherits \(\mathsf U_{\mathrm{fin}}(\mathcal O,\Omega)\) from \(X\). The relationship between \(X\) and \(Y\), the relevant topologies, and the apparent function-space membership error remain unresolved. The supplied registry labels this problem open; this draft does not independently verify that status.
\subsection{Sources}
\begin{itemize}
\item Local UnsolvedMath registry, permanent ID \(800622\), source code \texttt{catalogue entry}, title ``Selection Principles in Mathematics --- Problem 51 ([SFT]).'' The record attributes the problem to Tsaban, with proposed year 2006, and supplies the status label ``open.''
\item Supplied source attribution: Tsaban, \emph{Selection Principles and Special Sets of Reals: Open Problems} (2006), Problem 51. Supplied URL: \url{https://arxiv.org/abs/math/0609601}. The background records access on 2026-10-16 and extraction by \texttt{arxiv\_latex}; these are supplied provenance statements, not independently verified access or source verification.
\item The supplied attribution/context marker is \texttt{[SFT]}; its bibliographic expansion is absent from the record.
\end{itemize}
\end{document}
